\chapter{本课题创新之处}
\section{从年龄变量到辅助修正幅度的约束}
收益、文本、图和普通门控都已有研究基础。本课题拟检验的第一项差异，是在风险logit层明确区分主干和辅助修正，并把可用性、质量和年龄放入修正幅度的乘积包络。年龄只是普通输入时，预测头可以学习任意非线性关系，却未必满足资料陈旧时的最大作用限制；有界结构则将这一性质直接写进模型。

这种差异同时带来约束和代价。它使全缺失状态有明确回退，并提供随年龄缩小的允许偏移，但也会缩小函数类。有效旧信息可能因衰减而受到限制，解析质量标记也可能漏掉仍有信息的正文。研究贡献因此不是一个“必然提高”的不等式，而是检验结构限制在目标任务中的净效用。

\reporttable{年龄普通特征与有界修正的可检查差别}{tab:mechanismcomparison}{
\begin{tabularx}{\linewidth}{p{3.3cm}YY}\toprule
比较维度&年龄作为普通特征&年龄有界修正\\\midrule
资料陈旧度&由预测头任意利用&进入可学习非负衰减包络\\
最大辅助偏移&一般无显式年龄界&$M Q e^{-rA}$为分支上界\\
全缺失状态&依具体偏置/归一化实现&修正为0，回同模型主干\\
新鲜资料的利用&表达自由度较大&保留内容门和双向tanh修正\\
待测代价&噪声可能被放大&可能损害有效旧信息\\\bottomrule
\end{tabularx}}

表\ref{tab:mechanismcomparison}不是对所有门控方法的普遍定理。普通门控也可以被加入同样约束，因此真正需要比较的是具体实现、共同信息及编码器预算。提出的方法差异必须在同编码器控制、年龄仅作特征、去全部年龄和去质量实验中形成证据，而不是仅凭结构图认定创新。

\section{相对邻居权重与绝对陈旧度的两级处理}
归一化关系消息存在一个容易忽略的性质：目标同一名单的共同边龄因子可以抵消。仅写边龄权重，可能只改变不同邻居的相对贡献，却没有限制整份目标名单变旧时的整体影响。本课题将边级衰减和图外修正包络分开，前者处理邻居相对新旧，后者处理目标关系资料的绝对年龄。

去自身聚合进一步避免目标表示通过机构均值返回，独享机构不提供跨公司证据。目标自身仍可在图投影中与合法消息联合编码；关系随机化和同编码器预算用于检查增量是否来自容量，而非真实关系。此项属于面向本任务的消息设计和公平实现，不把GraphSAGE或去自身思想单独宣称为首创。

单一其他邻居、全部目标边同龄和高度机构是三个有价值的检验边界。前者使机构内权重抵消，第二种情况使外层年龄因子抵消，第三种情况产生资源与采样问题。这些边界使公式、数值检查和实证设计相互对应，也给“年龄究竟在哪一层起作用”提供具体解释。

\section{可追溯的工程合同与信息公平验证}
不同资料只要在名称和日期上能够拼接，并不意味着其经济对象和历史可用内容一致。本课题拟把来源、版本、主体、缓存和日索引贯穿到预测导出中，模型和强对照共享同一公司日。预告抽取与风险预测分开评价，正文信息同样给予强树，避免把输入差别当成融合增量。

这一工程合同的实践价值在于，结果不再只是一列概率。它能够说明某个分数使用了哪份公告、哪份名单、何种模型、哪些预处理参数，以及哪些资料因证据不足未进入。来源核验、全文语义评价和模型统计各自形成独立记录，在异常或版本修订时可以判断需要重做哪些步骤。

\section{创新主张需要满足的证据条件}
近期图时序和MSGCA已经覆盖与本方案接近的骨架，相关专利CN117808586A的公开线索也涉及风险目标和多关系融合。完整权利要求及部分近邻全文尚待核验，不能由摘要未出现某个词就推断方法不同。后续比较需逐项核对标签频率、资料可用性、关系重构、融合形式和评价协议。

若有界结构只相对收益-only模型提高，而不能超过相同信息普通门控或强树，则主要证据是新增信息而非融合创新。若增量只由单一年份或高覆盖群体驱动，应限定适用场景；若区间跨零，应报告未获得稳定增量。只有方法差异、性质、同信息比较和机制分析一致，才支持更强的研究贡献表述。投稿竞争力还涉及近邻充分性和真实结果，不由开题阶段预先保证。
